\chapter{Metodología}

\section{Enfoque y tipo de investigación}

Este trabajo se desarrolla bajo un enfoque de investigación aplicada y de tipo experimental. Es aplicada porque su resultado no es una contribución teórica independiente, sino un artefacto de software concreto y utilizable: la librería npm logSguarDian, con su pipeline de ingeniería de características y su modelo híbrido integrados en un middleware funcional para aplicaciones Express. Es experimental porque la validez de ese artefacto no se argumenta de forma discursiva, sino que se somete a criterios de aceptación predefinidos y cuantificados --los umbrales de F1-score, recall y tasa de falsos positivos establecidos en los objetivos específicos de este trabajo-- verificados empíricamente sobre un conjunto de prueba bloqueado y sobre un entorno de evaluación controlado y reproducible.

La combinación de ambos rasgos --construcción de artefacto y validación empírica contra criterios de aceptación predefinidos-- responde directamente al vacío metodológico identificado en el capítulo de Antecedentes. La literatura de referencia en detección de intrusiones mediante aprendizaje automático privilegia consistentemente los datasets de flujo de red heredados sobre la semántica de peticiones HTTP individuales. Ni las revisiones más recientes de arquitecturas híbridas supervisadas/no supervisadas ni las revisiones sistemáticas de cortafuegos de aplicación web potenciados por inteligencia artificial documentan, entre los estudios que cubren, un sistema implementado bajo una arquitectura embebida dentro del proceso de una aplicación Node.js con recursos limitados.

Al no existir un antecedente documentado que combine un modelo híbrido de aprendizaje automático con las restricciones de latencia, memoria y concurrencia de un solo hilo propias de un middleware Node.js, este trabajo no puede apoyarse en un diseño experimental replicado de la literatura. Requiere construir el artefacto y definir sus propios criterios de aceptación. Por ello, la clasificación que corresponde al problema es la de investigación aplicada-experimental, y no una alternativa entre varias igualmente válidas.

\section{Metodología de trabajo y fases}

El desarrollo del proyecto se organizó en seis fases secuenciales, cada una con una disciplina operativa asociada y un artefacto verificable como criterio de cierre (regla R3 del proyecto: ninguna fase se da por completa sin un número reproducible, una prueba que pasa, o un archivo committeado que la respalde).

\textbf{Fase 1 --- Construcción del pipeline de ingeniería de características (20 de junio -- 1 de agosto de 2026, objetivo específico 1).} Se implementó, en TypeScript, la extracción y vectorización de un vector de 76 características numéricas a partir del objeto \texttt{req} de Express, ejecutada de forma asíncrona mediante \textit{worker\_threads} para no bloquear el \textit{Event Loop} del proceso anfitrión (implementación inicial de 72 características a cargo de Sebastián Huertas; véase la sección~\ref{sec:recursos-participantes}). Esta fase está gobernada por la regla R1 del proyecto: la extracción de características vive únicamente en \texttt{packages/extractor} (TypeScript); ningún componente en Python recalcula features a partir de datos crudos, sino que consume matrices ya extraídas por el CLI del extractor. Esta disciplina se verificó de forma explícita incluso en la evaluación final del modelo (\texttt{training/evaluate\_test.py}), donde la selección de columnas se realiza por nombre contra la lista canónica \texttt{FEATURE\_NAMES} copiada literalmente del extractor TypeScript, en lugar de inferir el vector de entrada por descarte posicional. El artefacto de cierre de esta fase es el benchmark de latencia del extractor. Su p95 mixto con \textit{worker\_threads} es de 0.0273~ms (mediana de tres corridas), muy por debajo del límite de 5~ms establecido en el objetivo específico 1 (Cuadro \ref{tabla:res-extractor}).

\textbf{Fase 2 --- Construcción del corpus de entrenamiento.} Se ensambló un corpus a partir de nueve conjuntos auditados, ocho incluibles y siete incorporados, más dos listas públicas. Su licencia y procedencia se verificaron y se documentaron en la auditoría de datasets del proyecto. El corpus se complementó con un suplemento sintético acotado, generado mediante Docker Compose para cerrar vacíos de cobertura. Esos vacíos se detectaron mediante análisis de exclusión por fuente (\textit{Leave-One-Source-Out}) y pruebas de ablación sobre el corpus público. El suplemento nunca es la base primaria del corpus.

La fase cierra con la deduplicación exacta del corpus, el marcado de casi-duplicados y el bloqueo criptográfico del conjunto de prueba (\texttt{training/splits/test.lock.sha256}). Este bloqueo activa la regla R2: el conjunto de prueba se lee una sola vez para la evaluación final de cada generación de modelo. Ninguna decisión de ajuste posterior a esa lectura puede modificarse en función de los resultados observados en él. La generación v11 incumplió esta regla, como se detalla más adelante en este capítulo.

\textbf{Fase 3 --- Entrenamiento y calibración del modelo híbrido.} El entrenamiento de Random Forest e Isolation Forest se realizó íntegramente en Python sobre hardware Apple Silicon local, sin dependencia de infraestructura de cómputo en la nube. La selección de los umbrales de decisión se realizó sobre el conjunto de validación. En dos versiones de Isolation Forest esa selección se revisó después de leer el conjunto de prueba. En \texttt{if\_v2}, la confirmación arrojó una tasa de falsos positivos de 0.1011 frente a un límite de 0.10, y el objetivo de calibración pasó a FP $\leq$ 0.08. En \texttt{if\_v5}, la primera confirmación arrojó 0.0811 frente a un objetivo interno de 0.08, y el objetivo se endureció a FP $\leq$ 0.06.

En ambos casos el umbral final se seleccionó sobre validación, pero el conjunto de prueba influyó en el criterio de calibración. El umbral de \texttt{if\_v10} se calibró sobre validación como el máximo recall entre los umbrales que cumplen recall $\geq$ 0.50 y FP $\leq$ 0.06 (\texttt{retrain\_v11\_contract69.py}, líneas 107--110). El de Random Forest (\texttt{RF\_THRESHOLD} = 0.35) se calibró sobre validación para \texttt{rf\_v3} y se mantuvo sin recalibración para \texttt{rf\_v11}.

\textbf{Fase 4 --- Verificación de paridad ONNX.} Ambos modelos se serializaron a formato ONNX mediante \texttt{sklearn2onnx}, y su paridad de predicción frente al modelo original en \texttt{scikit-learn} \parencite{pedregosa2011} se verificó de forma automatizada (\texttt{training/models/parity\_report.json}). La verificación confirmó diferencias de probabilidad en el orden de $10^{-7}$ entre ambos entornos de inferencia. Esto ocurrió antes de habilitar el modelo serializado para su integración.

\textbf{Fase 5 --- Integración en el middleware.} El modelo serializado se integró en la librería npm mediante \texttt{onnxruntime-node}, ejecutado en una arquitectura de \textit{worker pool} dedicado (un worker para Random Forest y dos workers para Isolation Forest en hilos separados), diseño adoptado tras confirmar empíricamente que \texttt{onnxruntime-node} serializa las llamadas de inferencia concurrentes dentro de un mismo hilo nativo. Esta fase cierra con la evaluación del modelo híbrido sobre el conjunto de prueba bloqueado (objetivo específico 2, cierre 2 de agosto de 2026).

\textbf{Fase 6 --- Evaluación en tres configuraciones progresivas, la tercera con dos variantes (3a, solo cortafuegos; 3b, en capas) (2 de agosto -- 30 de septiembre de 2026, objetivo específico 3).} La librería integrada como middleware se evaluó en un entorno de prueba controlado (Docker Compose) bajo tres configuraciones sucesivas, la tercera con dos variantes (3a, solo cortafuegos; 3b, en capas): sin protección, con logSguarDian como única capa, y en defensa en profundidad junto con un Web Application Firewall perimetral. El refinamiento iterativo del modelo continuó dentro de esta fase. Incluyó el ciclo de reentrenamiento que atendió los vacíos de cobertura de la categoría de inyección de comandos detectados durante la evaluación. Ese refinamiento no alteró los criterios de aceptación establecidos en los objetivos específicos.

Los plazos de los objetivos específicos 2 y 3 (2 de agosto y 30 de septiembre de 2026) corresponden al protocolo. Las cifras finales de este trabajo se obtuvieron después de esas fechas. La lectura del conjunto de prueba de \texttt{rf\_v11} que se reporta es del 26 de septiembre. Las mediciones del coautor sobre la Ronda 5 constan en commits del 10 de octubre, y el paquete 0.1.1 se publicó el 11 de octubre de 2026 (UTC). La Ronda 5 con el par C se ejecutó el 9 de octubre. El refinamiento iterativo posterior al 2 de agosto se declara en el objetivo específico 2 y no alteró los criterios de aceptación.

\section{Técnicas e instrumentos}\label{sec:tecnicas-instrumentos}

La técnica central de la fase de construcción del modelo es la ingeniería de características sobre el objeto \texttt{req} de Express: un vector de 76 dimensiones organizado en grupos funcionales. Estos grupos capturan los marcadores semánticos de las cuatro categorías de ataque cubiertas (inyección SQL, \textit{cross-site scripting}, recorrido de directorios e inyección de comandos). Además, el vector incluye características estadísticas transversales (longitud, entropía de Shannon, proporción de caracteres especiales) y metadatos de la petición (método HTTP, encabezados, indicadores de codificación).

Conviene distinguir lo que el extractor produce de lo que cada modelo consume. De esas 76 dimensiones, el componente Random Forest consume 69 y el componente Isolation Forest consume 63. Cada modelo excluye las variables no disponibles en el momento de la intercepción de la petición (código de estado de respuesta, contadores de frecuencia por ventana temporal). Isolation Forest excluye además las variables con varianza próxima a cero sobre tráfico benigno.

En ambos modelos, de forma temporal, se excluye la característica más reciente del extractor (\texttt{non\_json\_quote\_count}, que descuenta las comillas de sintaxis JSON del conteo de comillas del payload). Esta última se excluye no por invalidez, sino porque su incorporación al vector de entrada de los modelos en producción implica un reentrenamiento completo. Esa promoción es una decisión pendiente y deliberadamente separada de su adición al extractor.

El Cuadro \ref{tabla:grupos-features} enumera los seis grupos funcionales y el número de variables que contiene cada uno, según la tabla de features de la especificación del extractor.

\begin{table}[htbp]
\centering
\footnotesize
\setlength{\tabcolsep}{3pt}
\begin{adjustbox}{max width=\linewidth}
\begin{tabular}{|c|p{5.6cm}|c|p{5.8cm}|}
\hline
\textbf{Grupo} & \textbf{Descripción} & \textbf{Variables} & \textbf{Posiciones en la especificación} \\ \hline
A & Semántica de ataque (SQLi, XSS, \textit{path traversal}, CmdI) & 37 & \#26--\#62 \\ \hline
B & Estructura HTTP & 9 & \#63--\#71 \\ \hline
C & URL y ruta & 6 & \#3--\#5, \#8--\#10 \\ \hline
D & Parámetros y cuerpo (composición del payload) & 9 & \#1, \#2, \#6, \#7, \#11--\#14, \#18 \\ \hline
E & Evasión por encoding y caracteres & 10 & \#15--\#17, \#19--\#25 \\ \hline
--- & Temporales (ventanas de 1, 5 y 60 segundos) & 5 & \#72--\#76 \\ \hline
\multicolumn{2}{|l|}{\textbf{Total}} & \textbf{76} & \\ \hline
\end{tabular}
\end{adjustbox}
\caption[Grupos de características del vector de 76 dimensiones]{Grupos de características del vector de 76 dimensiones, con el número de variables de cada uno.}
\label{tabla:grupos-features}
\par\smallskip\footnotesize Fuente: \texttt{docs/feature-spec.md} del repositorio logSguarDian, rama \texttt{docs/feature-spec-76}. Elaboración propia.
\end{table}

Las variables temporales (\#72--\#76) no son consumidas por ninguno de los dos modelos, porque requieren estado compartido entre peticiones que el extractor no mantiene por diseño. Tampoco lo son \texttt{status\_code} (\#71), que no existe al interceptar la petición, ni \texttt{non\_json\_quote\_count} (\#32). Random Forest consume así 69 variables; Isolation Forest consume 63, tras excluir además \texttt{null\_byte\_count}, \texttt{sensitive\_file\_target}, \texttt{dotdot\_encoded\_count}, \texttt{os\_path\_indicator}, \texttt{authorization\_length} y \texttt{unusual\_headers\_count}.

Sobre ese vector opera el modelo híbrido de aprendizaje automático descrito en el Marco Teórico: Random Forest \parencite{breiman2001} como componente supervisado, encargado de clasificar cada petición en una de las cuatro categorías de ataque cubiertas o como tráfico legítimo, e Isolation Forest \parencite{liu2008} como componente no supervisado, entrenado únicamente sobre tráfico benigno para aportar una señal complementaria de detección de anomalías estadísticas frente a variantes no vistas durante el entrenamiento. La justificación teórica de esta combinación --la complementariedad entre detección supervisada y no supervisada de una sola clase \parencite{chandola2009}, y la selección de métodos de ensamble ligeros frente a alternativas de redes neuronales de mayor capacidad-- se desarrolla en el Marco Teórico y no se repite aquí; esta sección se limita a documentar su aplicación instrumental. En su forma ONNX publicada, \texttt{rf.onnx} pesa 8.90~MB (8.49~MiB) e \texttt{if.onnx} 2.16~MB (2.06~MiB).

El instrumental de construcción y despliegue del modelo comprende \texttt{scikit-learn} \parencite{pedregosa2011} para el entrenamiento de ambos algoritmos, \texttt{sklearn2onnx} para su serialización a formato ONNX \parencite{bai2019}, \texttt{onnxruntime-node} como motor de inferencia embebido en el proceso Node.js, y Docker Compose. Docker Compose levanta tanto el entorno de generación de tráfico sintético (fase 2) como el entorno de prueba vulnerable de las tres configuraciones de evaluación (fase 6). La restricción de recursos de este instrumental --modelos de ensamble de complejidad computacional acotada, entrenados y servidos sin infraestructura de cómputo en la nube-- es consistente con la disciplina de inferencia embebida bajo restricciones de memoria y latencia documentada en la literatura de Edge AI \parencite{warden2019}.

Como técnica de interpretabilidad se empleó SHAP\parencite{lundberg2017shap}, específicamente su explicador \texttt{TreeExplainer} sobre el modelo Random Forest. Se aplicaron dos modalidades complementarias. La primera es la importancia global por clase, calculada sobre el conjunto de validación para identificar las diez características con mayor contribución promedio a cada categoría de ataque. La segunda es la explicación local de casos individuales, empleada para diagnosticar el comportamiento del modelo frente a peticiones específicas (por ejemplo, el efecto de la presencia o ausencia de encabezado \textit{User-Agent} sobre la puntuación de Isolation Forest). Los casos locales analizados mediante esta técnica se documentan con su evidencia numérica en el capítulo de Resultados.

El análisis se calculó sobre \texttt{rf\_v11.pkl} tras la corrección de trazabilidad de artefactos descrita en la sección de Recolección, procesamiento y validación de datos, por lo que las características y casos que explica corresponden al modelo efectivamente evaluado en este capítulo. La importancia global confirma que los encabezados de \textit{User-Agent} (\texttt{ua\_present}, \texttt{ua\_length}) figuran entre las características de mayor contribución. Esto ocurre no solo para la calibración de Isolation Forest ya documentada, sino también para la clasificación de Random Forest en las categorías \texttt{benign}, \texttt{sqli} y \texttt{path\_traversal}, con una contribución más débil, aunque presente, en \texttt{cmdi} y \texttt{xss}. Este predominio es desigual entre clases: es más marcado en tres de las cinco categorías y no debe leerse como uniforme.

Para \texttt{cmdi}, el análisis aporta evidencia cuantitativa adicional sobre la dificultad de separabilidad ya documentada. Incluso sobre el modelo corregido, las dos características de mayor contribución para esta clase (\texttt{method\_is\_get}, \texttt{ua\_length}) no son marcadores semánticos de inyección de comandos. Solo una característica semántica relacionada (\texttt{subshell\_count}) figura entre las diez de mayor peso. \texttt{shell\_command\_count} y \texttt{pipe\_count}, pese a formar parte del vector de entrada, no se encuentran entre las diez principales. Esto es consistente con, aunque no idéntico a, la dificultad estructural ya identificada entre \texttt{cmdi} y \texttt{path\_traversal}, y se documenta aquí como un hallazgo parcialmente resuelto, no cerrado.

\subsection{Pares de modelos evaluados}\label{sec:pares-modelos}

El paquete 0.1.0 y los modelos sobre los que se midieron las métricas de este trabajo no son el mismo par. Las métricas del capítulo de Resultados corresponden al par C, mientras que el paquete \path{logsguardian@0.1.0} del registro contiene el par A. El Cuadro \ref{tabla:pares-modelos} identifica cada par por el SHA-256 de sus archivos ONNX y registra en qué evaluaciones se usó.

\begin{table}[htbp]
\centering
\footnotesize
\setlength{\tabcolsep}{3pt}
\begin{adjustbox}{max width=\linewidth}
\begin{tabular}{|c|p{3.3cm}|p{2.9cm}|p{2.9cm}|p{4.4cm}|}
\hline
\textbf{Par} & \textbf{Modelos} & \textbf{SHA-256 (\texttt{rf.onnx} / \texttt{if.onnx})} & \textbf{Origen} & \textbf{Evaluaciones en que se usó} \\ \hline
A & \texttt{rf\_v10} / \texttt{if\_v9}, 67 y 61 características & \texttt{18083725\ldots} / \texttt{5de374a7\ldots} & Commit \texttt{dc2aaa65}, 2 de agosto de 2026 & Paquete 0.1.0 (par A), \path{logsguardian@0.1.0} (30 de agosto de 2026, commit \texttt{c3d591a}); Configuraciones 2, 3a y 3b sobre los 590 payloads (Ronda 4, 15 de agosto de 2026; sobre el solapamiento de ese corpus con el entrenamiento, véase la sección~\ref{sec:solapamiento-corpus})\textsuperscript{4}; Configuración 2 sobre 100 payloads por categoría (6 de agosto de 2026; el modelo consta en \path{docs/config2-latency-evaluation.md}, sin hash); Ronda 5 con el par A (tarball e09120\textellipsis{})\textsuperscript{3}; estudio de User-Agent de la aplicación vulnerable del proyecto (24 de septiembre)\textsuperscript{1} \\ \hline
B & \texttt{rf\_v11} / \texttt{if\_v10}, primera versión, 69 y 63 características & \texttt{a7c015f5\ldots} / \texttt{3cf8e45c\ldots} & Commit \texttt{734c24f}, 30 de agosto de 2026 & Primera lectura del conjunto de prueba (24 de septiembre); repetición de Config 2 sobre los 590 payloads (9 de septiembre)\textsuperscript{1}; experimento de latencia, memoria y CPU en Linux nativo \\ \hline
C & \texttt{rf\_v11} / \texttt{if\_v10}, reentrenamiento, 69 y 63 características & \texttt{25b407e6\ldots} / \texttt{5be28ee8\ldots} & Commit \texttt{03de50f}, 26 de septiembre de 2026 & Lecturas del 25 y 26 de septiembre del conjunto de prueba; paridad ONNX; análisis SHAP; Ronda 5 (9 de octubre de 2026) con el tarball de validación del commit \texttt{22550ea} (SHA-256 \texttt{cf1238\ldots}; véase la sección~\ref{sec:ronda5})\textsuperscript{3}; copia de Juice Shop\parencite{juiceshop}\textsuperscript{2} \\ \hline
\end{tabular}
\end{adjustbox}
\caption[Pares de modelos evaluados]{Pares de modelos evaluados, con su origen y las evaluaciones en que se usó cada uno.}
\label{tabla:pares-modelos}
\par\smallskip\footnotesize Fuente: commits \texttt{dc2aaa65}, \texttt{734c24f} y \texttt{03de50f} del repositorio logSguarDian; hashes de \texttt{training/models/} (\texttt{rf.onnx} e \texttt{if.onnx}); \texttt{docs/config3b-results.md} del repositorio logSguarDian-vulnerable-project. Elaboración propia.
\end{table}

\par\smallskip\footnotesize Notas: (1) El JSON de resultados de la repetición de Config 2 del 9 de septiembre y el del estudio de User-Agent no registran el modelo. Su atribución se apoya en el mensaje del commit \texttt{14c0e15} de la aplicación vulnerable del proyecto, que actualizó los modelos a \texttt{rf\_v11}/\texttt{if\_v10}, y en el contenido de \texttt{models/} en el commit del estudio (\texttt{5966712}). (2) De las aplicaciones de terceros DVNA\parencite{dvna}, node-api-goat\parencite{nodeapigoat} y dvws-node\parencite{dvwsnode} no consta el modelo exacto de cada medición; solo la medición de Juice Shop usa una copia cuyo modelo sí consta.

\footnotesize (3) La Ronda 5 con el par A (tarball e09120\textellipsis{}) se atribuye al par A según \texttt{docs/config3b-results.md} de la aplicación vulnerable del proyecto (rama main); el mismo documento atribuye el tarball del commit \texttt{22550ea} al par C. Los tres archivos se denominan \texttt{logsguardian-0.1.0.tgz}, pero son distintos. El tgz de la versión 0.1.0 del registro tiene SHA-256 \texttt{\seqsplit{c5c1a1da57cfb873350ac0c8ab3e7bad395e8cbf507b18cc3272df447880dcb0}} y contiene el par A; el de la versión 0.1.1 (11 de octubre de 2026, UTC) tiene SHA-256 \texttt{\seqsplit{abf8d5c47036d9c3a289bc36158f1caaa91b9504689aa97193e51c6808258d15}} y contiene el par C. El paquete 0.1.0 contiene el par A. El tarball de la Ronda 5 original (commit \texttt{a3aada0}, SHA-256 \texttt{\seqsplit{e09120006dd17e790838f2ae92bb68ee4c7f3a5e160c2c4c83f5a513909fe359}}) también contiene el par A, pero no es el tgz del registro. El tarball de validación y la copia local de Juice Shop, no publicados, contienen el par C.

\footnotesize Los nombres usados en este documento son cinco. El \textit{paquete 0.1.0 del registro} (\path{logsguardian@0.1.0}, par A; SHA-256 \texttt{\seqsplit{c5c1a1da57cfb873350ac0c8ab3e7bad395e8cbf507b18cc3272df447880dcb0}}). El \textit{tarball de la Ronda 5 original} (commit \texttt{a3aada0}, SHA-256 \texttt{\seqsplit{e09120006dd17e790838f2ae92bb68ee4c7f3a5e160c2c4c83f5a513909fe359}}, par A). El \textit{tarball de validación del commit 22550ea} (SHA-256 \texttt{\seqsplit{cf1238deb33383975e12058de71c4423b22d7c5bcba5cf47241c0d0f712b0416}}, par C). El \textit{paquete 0.1.1 del registro} (\path{logsguardian@0.1.1}, par C; SHA-256 \texttt{\seqsplit{abf8d5c47036d9c3a289bc36158f1caaa91b9504689aa97193e51c6808258d15}}). La \textit{copia local de Juice Shop} (SHA-256 \texttt{\seqsplit{c3497bddc69d68bb75eadf494d98b959b7dd4a2b2653cbe945fa2c242a5848a1}}, par C).

\footnotesize (4) La Configuración 1 no usa los modelos, porque la aplicación no tiene protección. En las Rondas 1 a 3, con corpus pequeño, el modelo no consta en los resultados.

Hashes completos. Par A: \texttt{rf.onnx} \texttt{\seqsplit{180837255dffd596a38bf7c29032ccc313005cd70eca5c6238e23e35bdc0d4ec}} e \texttt{if.onnx} \texttt{\seqsplit{5de374a7f1d19176d7c66d3354e972d37b28a7306a87efa0ce8639278f489150}}. Par B: \texttt{rf.onnx} \texttt{\seqsplit{a7c015f5a5fa3fabfd91b79ec9338dc7da6a68c392c19ded4b04cc6f0fb903b2}} e \texttt{if.onnx} \texttt{\seqsplit{3cf8e45c285de19bc86e6dc284e502cd2c24becb888747b9645bed4cfb224417}}. Par C: \texttt{rf.onnx} \texttt{\seqsplit{25b407e649283606a231665a5a8c8ea1c2c7ad4fd0374a8945b3f119651b0870}} e \texttt{if.onnx} \texttt{\seqsplit{5be28ee8437643b2f13708eee1ac3175348a350ea6fb8f8daeb09bd5717f797e}}.
% FUENTE: vendor/logsguardian-0.1.0.tgz del repositorio logSguarDian-vulnerable-project (commit 4d928b5); SHA-256 calculado con git show origin/main:vendor/logsguardian-0.1.0.tgz | shasum -a 256; sus models/rf.onnx e if.onnx son 25b407e6 y 5be28ee8.
% FUENTE: docs/evidence/README.md del repositorio logSguarDian-vulnerable-project (SHA-256 de rf.onnx e if.onnx calculados con git show del commit f15a9e8); la frase 'sin resolver' de docs/cybersecurity-objectives-compliance.md es anterior a esa verificación.

\subsection{Validación en aplicaciones de terceros}\label{sec:validacion-terceros}
% FUENTE: la opcionalidad de la validación con startup está en el protocolo aprobado y la recomendación del asesor no consta en el repositorio; el nombre del asesor está en 0-datos_estudiante.tex (\nombreasesor).
% FUENTE: modo de cada aplicación en docs/vulnerable-projects-tests/*-results.md y docs/limitations.md (addendum Juice Shop) de logSguarDian.

La validación con una startup guatemalteca, prevista en el protocolo, era opcional y no se realizó. Por recomendación del asesor, Ing. Erick Francisco Marroquín Rodríguez, el sistema se evaluó en cuatro aplicaciones vulnerables de terceros: DVNA, node-api-goat, dvws-node y OWASP Juice Shop. En las tres primeras se usó el middleware en modo \texttt{block} con la configuración por defecto (\texttt{RF\_THRESHOLD=0.35}), y se enviaron ataques reales junto con tráfico benigno.

En OWASP Juice Shop se usó el modo \texttt{monitor} y solo se enviaron peticiones benignas, porque con la tasa de falsos positivos observada no se generó corpus de ataque para esa aplicación. En esa misma aplicación se comparó el efecto del encabezado \textit{User-Agent}, con un valor de navegador y con \texttt{curl}. Los resultados se presentan en la sección~\ref{sec:falsos-positivos}.

\section{Recolección, procesamiento y validación de datos}

El corpus final reúne 383\,202 filas de 17 fuentes: 268\,240 en entrenamiento, 57\,481 en validación y 57\,481 en prueba. Siete provienen de conjuntos públicos auditados (CAPEC, modsec-learn, registros de ModSecurity \parencite{lucz2025}, Payloads.csv\parencite{shahane2021payloads}, payload\_full\parencite{thakur2025payloadfull}, XSS\_dataset\parencite{shah2020xss} y command injection\parencite{pawase2024cmdi}). Dos provienen de listas públicas de payloads (SecLists\parencite{seclists} y PayloadsAllTheThings\parencite{payloadsallthethings}, ambas con licencia MIT). Siete provienen de generadores propios y una de telemetría curada del pipeline de MLOps (119 filas de una ejecución simulada, no de tráfico real; anexo~\ref{anexo:loso}).

La auditoría evaluó nueve conjuntos: ocho se consideraron incluibles y uno se excluyó por incompatibilidad de taxonomía de etiquetas (AIT-LDSv2, porque sus etiquetas describen roles del atacante y no tipos de payload). Uno de los ocho, omurugur, no se incorporó al corpus final. Las licencias de los conjuntos auditados se verificaron con las API de Zenodo, Harvard Dataverse y GitHub y con los metadatos de Kaggle; dos conjuntos quedan con licencia desconocida (\texttt{XSS\_dataset} y \texttt{command injection}).

Ese suplemento sintético no es una fuente primaria del corpus. Constituye un conjunto de correcciones dirigidas. \texttt{synthetic\_nav} y \texttt{synthetic\_nav\_ecommerce} se generaron para cerrar un vacío de generalización sobre tráfico benigno de navegación pura detectado mediante \textit{Leave-One-Source-Out}. Al excluir \texttt{synthetic\_nav} del entrenamiento, el F1 macro sobre sus filas fue de 0.0027, el menor de las nueve fuentes evaluadas según \texttt{docs/limitations.md}. Tras añadir \texttt{synthetic\_nav\_ecommerce} al corpus, fue de 0.2916 (\texttt{loso\_results.json}; el archivo no registra el modelo, y el script entrena un Random Forest nuevo con los hiperparámetros de \texttt{rf\_v7}, \texttt{n\_estimators=30} y \texttt{max\_depth=25}). El valor 0.0027 consta solo como texto de una nota de ese archivo, no como métrica.

\texttt{synthetic\_multifield} se generó para reforzar la señal de formularios multi-campo tras el hallazgo del bug de desempate en \texttt{extractBestPayload}. Por su parte, \texttt{synthetic\_cmdi\_minimal\_context}, \texttt{synthetic\_cmdi\_windows} y \texttt{synthetic\_cmdi\_compound} se generaron específicamente para cerrar los vacíos de separabilidad de la categoría de inyección de comandos identificados mediante pruebas de ablación. Su justificación completa se describe en el Marco Teórico y en la sección de limitaciones documentadas del proyecto. En ningún caso el suplemento sintético sustituye al corpus público como base del entrenamiento: cada adición está motivada por un hallazgo específico de cobertura, no por una decisión general de completar el corpus con datos generados.

La validación del corpus antes de su uso comprendió dos procedimientos concretos. El primero es la verificación de licencias: en los dos casos donde el propio proveedor (Kaggle) reporta la licencia como no especificada (\texttt{XSS\_dataset.csv} y \texttt{command injection.csv}), esa ausencia se documenta explícitamente como tal en lugar de asumir una licencia permisiva por defecto. El segundo corresponde a la deduplicación del pipeline de unificación (\texttt{training/unify.py}) y a una auditoría aparte de casi-duplicados, separadas por su alcance.

El pipeline elimina las filas con huella exacta idéntica (path, query, body y etiqueta; función \texttt{deduplicate\_exact}, línea 94) y solo marca, sin eliminar, los pares de consultas con distancia de Levenshtein menor que 3 dentro de la misma clase (líneas 179--181; campo \texttt{query} de 1 a 100 caracteres, con tope de 10\,000 comparaciones por clase). La auditoría usó MinHash\parencite{broder1997} con \textit{banding} de LSH ajustado explícitamente (\texttt{num\_perm=128}, \texttt{params=(8,16)}), frente a la configuración por defecto derivada del umbral, que resultaba intratable en tiempo de cómputo sobre las clases de mayor volumen (\texttt{docs/limitations.md}, commit \texttt{bab5ce6} de logSguarDian).

Esa medición confirmó una saturación de duplicados considerable en las cinco clases: la clase \texttt{sqli} alcanza aproximadamente 395 veces más pares casi-duplicados que filas (dominada por generación por plantilla de \texttt{sqlmap}/CAPEC), y la clase \texttt{benign}, aproximadamente 9.5 veces, consistente con tráfico de cola larga y no con contenido generado por plantilla. La decisión tomada es el marcado sin eliminación de los casi-duplicados, porque remover filas reconfigura el balance de clases y requiere su propia estrategia de reentrenamiento. La cuantificación por partición se presenta en el anexo~\ref{anexo:casi-duplicados}.

% Fuente de la identidad de predicciones sobre validación: training/audit/identity_audit.json, rama chore/audit-identity-script
% de logSguarDian (onnx_comparison_on_val: argmax_differences 0, max_abs_probability_difference 0.0, 57 481 filas). Compara el
% rf.onnx validado (sha256 25b407e6…) con el rf_current_split.onnx (sha256 a4dddbb7…); el audit no compara las lecturas del 25 y
% el 26 de septiembre entre sí.
El procesamiento del corpus hacia los conjuntos de entrenamiento, validación y prueba sigue un procedimiento secuencial y verificable, no una descripción abstracta de buenas prácticas. Primero, se unifican las fuentes dispares en un esquema común de columnas. Segundo, se extrae el vector de 76 características sobre cada registro unificado (de las cuales 69 alimentan a Random Forest y 63 a Isolation Forest, según se detalla en la sección de Técnicas e instrumentos). Esa extracción usa el mismo CLI del extractor TypeScript que se emplea en producción (regla R1: nunca se recalculan features en Python).

Tercero, se realiza la partición en conjuntos de entrenamiento, validación y prueba. Cuarto, se aplica el bloqueo criptográfico del conjunto de prueba. El hash SHA-256 del bloqueo se guarda en \texttt{training/splits/test.lock.sha256}, antes del reentrenamiento de v3, v4, v6, v8, v9, v10 y v11, según los mensajes de los commits. El lock de v11 se re-fijó el 23 de septiembre de 2026, después de su reentrenamiento (anexo~\ref{anexo:locks}). Quinto, el entrenamiento y la calibración de ambos modelos se realizan exclusivamente contra los conjuntos de entrenamiento y validación.

La regla R2 --el conjunto de prueba se lee exactamente una vez para la evaluación final de cada generación de modelo-- no es una declaración de intención. Es un procedimiento verificable en el propio repositorio. La evaluación final de \texttt{rf\_v11}/\texttt{if\_v10} (\texttt{training/evaluate\_test.py}) recalcula el hash del conjunto de prueba restaurado y lo compara contra el valor bloqueado antes de ejecutar cualquier inferencia. Si los valores no coinciden, la evaluación aborta. El resultado de esa verificación queda registrado junto con las métricas en \texttt{training/results/v11\_test\_results.json} (\texttt{lock\_verified: true}).

Para la generación v11, el conjunto de prueba se leyó en tres ocasiones, por lo que el criterio de lectura única de R2 no se cumplió. El 24 de septiembre de 2026 se evaluó el modelo entrenado el 30 de agosto (macro F1 de 0.9843) sobre el conjunto re-bloqueado el 23 de septiembre; el repositorio no conserva la huella de la partición con que ese modelo se entrenó, por lo que no puede descartarse solapamiento. El 25 de septiembre se evaluó un modelo reentrenado sobre la partición vigente, como parte de la comparación con MLP y autoencoder. El 26 de septiembre se evaluó el modelo vigente en el repositorio, reconstruido con el mismo procedimiento porque el artefacto serializado original había sido sobrescrito. Estas dos últimas lecturas coinciden en macro F1 (0.9776) y sus predicciones sobre validación son idénticas.

Las métricas de este capítulo corresponden a la lectura del 26 de septiembre. El historial no registra ninguna decisión de modelado atribuida al resultado de la primera lectura. El umbral de Isolation Forest de la primera lectura (0.00807) corresponde al modelo del 30 de agosto; el de las dos siguientes (0.00421), al modelo reentrenado, cuyo umbral se calculó sobre validación. El historial no atribuye ese cambio al resultado de la primera lectura, pero no puede descartarse su influencia.

La confirmación en prueba del umbral de 0.35 con \texttt{rf\_v3}, afirmada en el commit \texttt{e74a8c6}, no tiene artefacto de respaldo en el repositorio. Dos limitaciones adicionales: 261 de las 1\,468 filas de \texttt{cmdi} del conjunto de prueba (17.8\,\%) son sintéticas y su generador no está versionado, y la fuente de los datos restaurados el 23 de septiembre no está documentada.

Durante la preparación de este capítulo se detectó y corrigió una brecha de trazabilidad de artefactos en el pipeline de evaluación, documentada aquí por transparencia metodológica. El archivo \texttt{training/models/rf\_v11.pkl}, al no estar bajo control de versiones, fue sobrescrito silenciosamente por una versión candidata --generada durante una investigación no relacionada sobre diversificación de tráfico benigno sintético, y explícitamente no promovida a producción-- bajo el mismo nombre de archivo. Esto dejó huérfano el emparejamiento exacto entre modelo y conjunto de prueba que una evaluación previamente citada había medido. El mecanismo de bloqueo criptográfico del conjunto de prueba en sí no se violó. Ese bloqueo ancla la identidad de cada fila, no el valor de sus características, por lo que es correctamente invariante a cambios del extractor.

Sin embargo, ninguna suma de verificación vinculaba una métrica citada con el contenido exacto, byte a byte, del par modelo/conjunto de prueba que la produjo. Se reentrenaron \texttt{rf\_v11}/\texttt{if\_v10} contra el contrato de 69/63 features vigente en producción (\texttt{parity\_report.json}) y se evaluaron sobre el conjunto de prueba correctamente re-bloqueado. Las cifras reportadas en este capítulo corresponden a la lectura del 26 de septiembre de 2026. El mismo conjunto se había leído también el 24 de septiembre (commit \texttt{3086318}) y el 25 de septiembre de 2026 (commit \texttt{fc47a5b}).

Como salvaguarda estructural se incorporaron tres mecanismos. El primero es un manifiesto de sumas de verificación sha256 (conjunto de prueba y ambos modelos ONNX) que acompaña ahora cada evaluación citada. El segundo es una extensión del pipeline de integración continua de entrenamiento que respalda y restaura \texttt{training/splits/} de la misma forma en que ya protegía \texttt{training/models/}. El tercero es el paso de \texttt{rf\_v11.pkl}/\texttt{rf\_v11\_metadata.json} a control de versiones, para que el modelo efectivamente evaluado no pueda volver a sobrescribirse sin dejar rastro.

\section{Métricas e indicadores}

Ninguna métrica de este trabajo se reporta como \textit{accuracy} global. El corpus es fuertemente desbalanceado --la clase \texttt{benign} y la clase \texttt{sqli} concentran, entre las dos, más de la mitad de los registros, mientras que \texttt{cmdi} es la categoría de ataque con menor representación-- y bajo ese desbalance un clasificador que prediga siempre la clase mayoritaria obtendría una \textit{accuracy} alta sin haber aprendido a distinguir ninguna categoría de ataque. Por esa razón, cada componente del sistema se evalúa con la métrica que corresponde a su naturaleza --clasificación multiclase supervisada en un caso, detección de anomalías de una sola clase en el otro-- y no con una métrica agregada que oculte ese desbalance.

Para Random Forest, la métrica de aceptación es el F1-score por clase: la media armónica de precisión y recall, que penaliza específicamente el sacrificio de una a expensas de la otra, a diferencia de reportar ambas por separado o de promediarlas aritméticamente. Esto importa de forma directa para este sistema: un modelo que maximizara solo precisión bloquearía pocas peticiones legítimas pero dejaría pasar ataques (bajo recall), y uno que maximizara solo recall bloquearía tráfico legítimo con frecuencia inaceptable para un middleware de producción (baja precisión); el F1-score es la métrica que exige un balance entre ambos costos operativos. El objetivo específico 2 fija el criterio de aceptación en F1-score mínimo de 0.80 en al menos tres de las cuatro categorías de ataque cubiertas, y no en las cuatro.

Esto reconoce explícitamente la dificultad estructural de separabilidad de \texttt{cmdi} documentada en el Marco Teórico y en las limitaciones del proyecto. La sintaxis de inyección de comandos comparte marcadores léxicos (rutas de sistema de archivos, operadores de redirección) con la categoría de recorrido de directorios. Esto reduce estructuralmente su separabilidad frente a las otras tres categorías, cuyos marcadores semánticos son más distintivos entre sí. El umbral de tres de cuatro clases no es, por tanto, una relajación arbitraria del criterio, sino el reconocimiento de una dificultad de separabilidad ya documentada antes de fijar el criterio de aceptación.

Sobre el conjunto de prueba bloqueado, la generación de modelos vigente (\texttt{rf\_v11}/\texttt{if\_v10}) alcanza un F1-score macro de 0.9776, con F1 por clase de 0.9984 (\texttt{benign}), 0.9239 (\texttt{cmdi}, precisión 0.9015 / recall 0.9475), 0.9826 (\texttt{path\_traversal}), 0.9956 (\texttt{sqli}) y 0.9877 (\texttt{xss}). Las cinco clases superan el umbral de 0.80 del objetivo específico 2, incluida la clase estructuralmente más difícil. Como evidencia complementaria se calculó también el área bajo la curva ROC uno-contra-el-resto \parencite{fawcett2006} para cada clase, con valores entre 0.99899 y 0.99998, que indica que, dentro del corpus, la separación entre clases no depende del punto de corte de probabilidad elegido para la decisión de bloqueo. Las métricas por clase se calculan sobre la clase de mayor probabilidad (argmax), sin el umbral de confianza de 0.35 con el que el sistema decide el bloqueo; no se calculó una métrica de prueba con la regla de decisión desplegada.

Para Isolation Forest, la métrica de aceptación es distinta por diseño: recall de detección de anomalías y tasa de falsos positivos sobre tráfico benigno, no F1-score. La razón es la naturaleza no supervisada del componente: Isolation Forest se entrena únicamente sobre tráfico benigno y no tiene conocimiento de las categorías de ataque como clases distintas, por lo que no existe una matriz de confusión multiclase sobre la cual calcular precisión o F1 por categoría; su tarea es binaria (anómalo o no) y se evalúa como tal. El protocolo original fija el criterio de aceptación en recall de detección mínimo del 50\% con una tasa de falsos positivos máxima del 10\%. Ambos criterios se cumplen sobre el conjunto de prueba bloqueado (\texttt{training/results/v11\_test\_results.json}, lectura verificada contra el bloqueo criptográfico, correspondiente a la evaluación del 26 de septiembre de 2026 descrita en este capítulo) con la generación vigente. Se obtuvo recall agregado de 0.9125 y tasa de falsos positivos de 0.0546.

Ese umbral del 10\% de falsos positivos se endureció operativamente durante iteraciones previas del proyecto (generación \texttt{if\_v5}) a un objetivo de calibración de 6\% en el conjunto de validación. No fue una preferencia estética, sino una lección concreta de margen insuficiente. Una calibración anterior de \texttt{if\_v5} (umbral 0.02868, con un objetivo interno de calibración de 8\% en validación) dio en el conjunto de prueba una tasa de falsos positivos de 0.0811. Es decir, quedó 0.0011 por encima de ese objetivo interno.

El límite formal del objetivo específico 2 se mantenía en 10\%. El corrimiento observado entre validación y prueba se atribuyó al cambio de features introducido junto con \texttt{rf\_v6} (el análisis por campo del cuerpo de la petición, que redujo la señal de longitud de payload de la que Isolation Forest dependía parcialmente). Ese cambio se introdujo antes de la calibración de \texttt{if\_v5} dentro del mismo ciclo de reentrenamiento, y motivó la recalibración.

El modelo recalibrado de \texttt{if\_v5} (umbral 0.00941) no es el de producción: el modelo vigente, \texttt{if\_v10}, usa el umbral 0.00421 y también se calibró en validación contra una tasa de falsos positivos $\leq$ 6\%. Desde entonces, la calibración de cada nueva generación de Isolation Forest elige como umbral el de máximo recall entre los que cumplen recall $\geq$ 0.50 y FP $\leq$ 6\% (\texttt{retrain\_v11\_contract69.py}, líneas 107--110). Así se deja margen deliberado de holgura frente al límite duro de 10\% fijado por el objetivo específico 2.

La verificación de paridad ONNX se reporta como una métrica de validación en sí misma, no como un detalle de implementación: dado que el entrenamiento ocurre en Python sobre \texttt{scikit-learn} y la inferencia en producción ocurre en Node.js sobre \texttt{onnxruntime-node}, cualquier discontinuidad entre ambos entornos invalidaría todas las métricas anteriores como descripción del comportamiento real del sistema en producción. La paridad medida entre ambos entornos es del orden de $10^{-7}$ en la diferencia de probabilidad de salida, confirmada de forma automatizada antes de habilitar cada modelo para su integración (\texttt{training/models/parity\_report.json}, \texttt{parity\_passed: true}). Esta métrica es la que permite afirmar que el F1-score, el recall y la tasa de falsos positivos reportados para el modelo en Python son también válidos como descripción del modelo efectivamente servido por el middleware.

Finalmente, la cobertura de detección por categoría en las configuraciones de evaluación integradas (Config 2 y Config 3) es una métrica distinta de las anteriores por su nivel de análisis. Mientras el F1-score y el recall/FP de Isolation Forest evalúan el modelo de forma aislada sobre un conjunto de prueba estático, la cobertura por categoría evalúa el sistema integrado --extractor, modelos e inferencia en \textit{worker\_threads}-- contra tráfico HTTP real dirigido a una aplicación en ejecución. El objetivo específico 3 fija este criterio exactamente como sigue: ``verificando que la detección de payloads representativos alcance al menos el 80\% por cada una de las cuatro categorías de ataque cubiertas en un entorno de prueba controlado, sin elevar la latencia p95 del servidor más de 15~ms respecto a la línea base sin protección''. Sobre la Ronda 5, con solapamiento bajo con el entrenamiento (21 de 200 peticiones de \texttt{cmdi} con un vecino con $J \geq 0.8$; ninguna en \texttt{sqli} ni en \texttt{path\_traversal}; anexo~\ref{anexo:casi-duplicados}), el umbral de 80\% se cumple con \texttt{rf\_v11}/\texttt{if\_v10} en \texttt{sqli} (99.5\%), \texttt{path\_traversal} (95.0\%) y \texttt{cmdi} (96.0\%) (Cuadro~\ref{tabla:res-ronda5}). La categoría \texttt{xss} carece de corpus limpio, y su cobertura de 97.3\% proviene del corpus de 590 payloads de SecLists (Cuadro~\ref{tabla:res-config3-corpus}). Los 100 payloads por categoría de Config 2 (Cuadro~\ref{tabla:res-config2-cobertura}) coinciden con datos de entrenamiento (sección~\ref{sec:solapamiento-corpus}) y no se usan como evidencia de cumplimiento.

Se entiende por payloads representativos los 200 de cada categoría de la Ronda 5, para \texttt{sqli}, \texttt{path\_traversal} y \texttt{cmdi} (Cuadro~\ref{tabla:res-ronda5}), y, para \texttt{xss}, los del corpus de 590 payloads de SecLists (Cuadro~\ref{tabla:res-config3-corpus}), porque la Ronda 5 no incluye esa categoría. El criterio de 80\% se evalúa sobre el valor observado en la muestra. El intervalo de Wilson\parencite{wilson1927} al 95\% se reporta como medida de incertidumbre y no como regla de decisión. Cuando su límite inferior queda bajo 80\%, el cumplimiento se califica como frágil.

Los criterios de memoria y de CPU provienen del protocolo del trabajo de graduación y no forman parte de los objetivos específicos. Se asocian al requisito del objetivo general de no comprometer el rendimiento del servidor anfitrión. El criterio de memoria fija un incremento máximo de 150~MB (unos 143~MiB) atribuible al middleware, no a la memoria absoluta del contenedor. El criterio de CPU fija un incremento relativo máximo de 15\,\% en carga normal y de ataque. El protocolo fijaba también un incremento de hasta 10\,\% de la latencia media, que no se midió (sección IX.5). Los resultados se reportan en la sección~\ref{sec:memoria-cpu}.

El criterio de 5~ms del objetivo específico 1 se aplica al p95 del extractor en aislamiento, con ida y vuelta a \textit{worker\_threads} (Cuadro~\ref{tabla:res-extractor}). El sobrecosto del middleware completo se evalúa en el objetivo específico 3.

El objetivo general no define una métrica propia. Se evalúa mediante los tres objetivos específicos, los criterios de rendimiento (latencia, memoria y CPU) y, para las variantes desconocidas, la marca de anomalía de Isolation Forest. Para esta última no se fijó un criterio de aceptación, por lo que solo se reporta evidencia indirecta (Cuadro~\ref{tabla:disc-cumplimiento}). La tasa de bloqueo de tráfico benigno en aplicaciones de terceros tampoco tiene criterio de aceptación en el protocolo, y se reporta como evidencia complementaria sobre el desempeño fuera del corpus, no como criterio de cumplimiento.

Como indicador complementario --de interpretabilidad, no de desempeño-- se reporta también el análisis SHAP descrito en la sección anterior. La importancia global por clase documenta qué grupos de características sostienen cada F1-score por categoría. Las explicaciones locales documentan casos específicos donde el modelo se comporta de forma distinta a lo esperado (por ejemplo, el efecto del encabezado \textit{User-Agent} sobre la puntuación de Isolation Forest). Esta evidencia complementa las métricas agregadas anteriores sin sustituirlas. Los hallazgos de esta importancia global, incluido el carácter parcialmente resuelto de \texttt{cmdi}, se detallan en la sección de Técnicas e instrumentos.

\section{Matriz de trazabilidad}

El Cuadro \ref{tabla:trazabilidad} articula la correspondencia entre cada objetivo del proyecto, la actividad concreta de la fase de trabajo que lo atiende, el instrumento o técnica empleada, la métrica de aceptación aplicada y el resultado real medido, con su fuente verificable en el repositorio del proyecto.

{\setlength{\tabcolsep}{3pt}
\begin{longtable}{|p{2.1cm}|p{2.6cm}|p{2.6cm}|p{2.5cm}|p{4.4cm}|}
\hline
\textbf{Objetivo} & \textbf{Actividad} & \textbf{Instrumento/\allowbreak técnica} & \textbf{Métrica} & \textbf{Resultado real} \\ \hline
\endfirsthead
\hline
\textbf{Objetivo} & \textbf{Actividad} & \textbf{Instrumento/\allowbreak técnica} & \textbf{Métrica} & \textbf{Resultado real} \\ \hline
\endhead

Específico 1 (pipeline de features, 20 jun -- 1 ago 2026) & Fase 1: construcción del pipeline de ingeniería de características & Extractor TypeScript + \textit{worker\_threads} (regla R1) & Latencia p95 de extracción $\leq$ 5~ms por petición & p95 mixto = 0.0273~ms con \textit{worker\_threads} (microbenchmark en aislamiento, mediana de tres corridas; Cuadro \ref{tabla:res-extractor}); diferencia no distinguible de 0 en WSL2 con resolución de 1~ms (compilación de julio de 2026, Cuadro~\ref{tabla:res-overhead}) \\ \hline

Específico 2 --- RF (modelo híbrido, cierre 2 ago 2026) & Fase 3: entrenamiento y calibración; Fase 4: verificación de paridad ONNX & Random Forest (\texttt{scikit-learn} + ONNX) sobre conjunto de prueba bloqueado (R2) & F1-score $\geq$ 0.80 en al menos 3 de 4 categorías de ataque & F1 macro = 0.9776; por clase: \texttt{benign} 0.9984, \texttt{cmdi} 0.9239, \texttt{path\_traversal} 0.9826, \texttt{sqli} 0.9956, \texttt{xss} 0.9877 --- las 5 clases superan el umbral \\ \hline

Específico 2 --- IF (modelo híbrido, cierre 2 ago 2026) & Fase 3: entrenamiento y calibración; Fase 4: verificación de paridad ONNX & Isolation Forest (\texttt{scikit-learn} + ONNX) sobre conjunto de prueba bloqueado (R2) & Recall de anomalías $\geq$ 50\%, tasa de falsos positivos $\leq$ 10\% & Recall = 0.9125, FP = 0.0546 --- ambos criterios cumplidos con margen \\ \hline

Específico 2 --- paridad (prerrequisito de integración) & Fase 4: verificación de paridad ONNX & Comparación de salida \texttt{scikit-learn} (Python) vs. \texttt{onnxruntime-node} & Diferencia de probabilidad entre ambos entornos & $\sim 10^{-7}$; \texttt{parity\_passed: true} (\texttt{parity\_report.json}) \\ \hline

Específico 3 --- cobertura (2 ago -- 30 sep 2026) & Fase 6: evaluación en Config 2 (+logSguarDian) & Middleware integrado sobre app vulnerable en Docker Compose & Detección $\geq$ 80\% de payloads representativos por categoría & Ronda 5, sin solapamiento con el entrenamiento, 200 payloads por categoría, par A (tarball e09120\textellipsis{}) y par C validado: \texttt{sqli} 100\,\% y 99.5\,\%, \texttt{path\_traversal} 85.0\,\% y 95.0\,\%, \texttt{cmdi} 91.5\,\% y 96.0\,\% (secciones \ref{sec:solapamiento-corpus} y \ref{sec:ronda5}). \texttt{xss} no tiene corpus limpio y no se evaluó en la Ronda 5. Config 2, 100 payloads por categoría: \texttt{sqli} 100\%, \texttt{xss} 100\%, \texttt{path\_traversal} 93\%, \texttt{cmdi} 100\%, con solapamiento con el entrenamiento (sección~\ref{sec:solapamiento-corpus}) \\ \hline

Específico 3 --- latencia (2 ago -- 30 sep 2026) & Fase 6: evaluación en Config 2 (+logSguarDian) & Benchmark de latencia (Artillery\parencite{artillery}; WSL2, Linux nativo y Docker+Postgres) & $\Delta$p95 $\leq$ 15~ms respecto a la línea base sin protección & diferencia no distinguible de 0 en WSL2 (resolución de 1~ms) y de 7--9~ms en Docker+Postgres, dentro del límite. En Linux nativo (GitHub Actions, ejecución 36276690236, 2 vCPU), $\Delta$p95 de 12.08~ms en normal y 12.85~ms en ataque, dentro del límite, y de 6\,495.5~ms en volumen con 7.00\,\% de timeouts, que no cumple (sección~\ref{sec:config2}) \\ \hline

Específico 3 --- defensa en profundidad (Config 3) & Fase 6 extendida: evaluación en Config 3a/3b (WAF solo / WAF + logSguarDian) sobre un corpus de 590 payloads de SecLists (solapamiento con el entrenamiento en la sección \ref{sec:solapamiento-corpus}) & ModSecurity v3 + OWASP CRS (PL1) como capa perimetral, logSguarDian como segunda capa & Cobertura agregada por configuración; capacidad de logSguarDian de detectar de forma independiente los ataques que evaden el WAF & Config 3a (solo WAF): 550/590 (93.2\%), con brechas reales por categoría (\texttt{sqli} 85.7\%, \texttt{path\_traversal} 87.0\%). Config 3b (capas): 587/590 (99.5\%). De los 40 ataques que evadieron el WAF, 37 fueron capturados de forma independiente por logSguarDian (92.5\% del vacío cerrado) \\ \hline

General (sin comprometer el rendimiento) & Fase 6: evaluación de Config 2 en Linux nativo & Contenedor completo medido con \texttt{docker stats}, 2 vCPU & $\Delta$Memoria $\leq$ 150~MB ($\approx$ 143~MiB); $\Delta$CPU relativo $\leq$ 15\,\% & +331.9, +332.0 y +375.9~MiB (ataque, normal y volumen); $\Delta$CPU de +529.6\,\%, +614.4\,\% y +40.1\,\%; no cumple (sección~\ref{sec:memoria-cpu}) \\ \hline

\caption[Matriz de trazabilidad del proyecto]{Correspondencia entre objetivos, actividades, instrumentos, métricas y resultados verificados. Elaboración propia. Las fechas de cierre de los objetivos son las del protocolo. La fecha de cada medición final se indica en la sección VII.2.}
\label{tabla:trazabilidad}
\end{longtable}}

\section{Recursos, participantes y restricciones}\label{sec:recursos-participantes}

En cuanto a recursos, el entrenamiento de los modelos se realizó sobre hardware local (Apple Silicon), sin presupuesto asignado a instancias de cómputo en la nube --condición ya establecida como restricción de recursos en el capítulo de Alcance y verificada como viable dado el tamaño acotado de los modelos de ensamble seleccionados frente a alternativas de mayor capacidad computacional--. El entorno de evaluación integrado (Config 1/2/3) se ejecutó igualmente sobre infraestructura local mediante Docker Compose, sin servicios de terceros ni infraestructura de cómputo remota.

En cuanto a participantes, este documento corresponde al trabajo de graduación individual de Diego Pablo Valenzuela Palacios. Documenta específicamente el componente de aprendizaje automático de logSguarDian: el pipeline de ingeniería de características, el modelo híbrido Random Forest/Isolation Forest, su serialización y verificación de paridad ONNX, y su evaluación integrada como middleware. El alcance exacto queda delimitado en el capítulo de Alcance de este trabajo.

El desarrollo del proyecto logSguarDian en su conjunto contó además con la colaboración de Sebastián Huertas en el componente de ciberseguridad y en la implementación del lado npm de la librería (extractor TypeScript y paquete \texttt{core}). Esa colaboración se reconoce aquí por transparencia metodológica. No implica coautoría de este documento ni extiende su alcance evaluado más allá del componente de Machine Learning aquí documentado. El extractor inicial de 72 características y el paso a \textit{worker\_threads} los implementó Sebastián Huertas. Diego Valenzuela analizó y justificó cada característica, definió cuáles no existen en inferencia y fijó el contrato de entrada, añadió cuatro características y las correcciones del extractor surgidas del análisis del modelo, y desarrolló el entrenamiento, la serialización ONNX y el pipeline de reentrenamiento.

En cuanto a restricciones, el proyecto se ejecutó dentro de la ventana comprendida entre mayo y octubre de 2026, conforme al calendario del trabajo de graduación de la Facultad de Ingeniería de la Universidad del Valle de Guatemala. Cualquier línea de trabajo no completable dentro de dicho período se documenta explícitamente como trabajo futuro y no como un entregable pendiente sin reconocer. La configuración de evaluación en defensa en profundidad (Config 3) completó su ronda de mayor escala (590 payloads de SecLists, cuyo solapamiento con el entrenamiento se describe en la sección \ref{sec:solapamiento-corpus}) dentro de esa ventana. El resultado se reporta en el Cuadro \ref{tabla:trazabilidad}. El corpus expuso brechas reales del Web Application Firewall que opera solo (93.2\% de cobertura agregada, con categorías tan bajas como 85.7\%), y logSguarDian, que opera como segunda capa, capturó 37 de los 40 ataques que evadieron esa primera capa.

Esa misma ronda dejó, no obstante, un hallazgo de calibración que motivó el estudio de falsos positivos descrito en las secciones~\ref{sec:config2} y~\ref{sec:falsos-positivos}. Una petición benigna de bajo contenido semántico puede acercarse lo suficiente al umbral de decisión de Random Forest. En ese caso, un metadato incidental de la petición (la cadena de \textit{User-Agent} del cliente HTTP, no el contenido del payload) determina el veredicto. Por ello se realizó un estudio de tasa de falsos positivos contra tráfico legítimo diverso, independiente de esta comparación basada en corpus de ataque.

Ampliar el corpus benigno queda como trabajo futuro, según se indica en el capítulo de Recomendaciones (\ref{cap:recomendaciones}). De igual manera, y como ya se estableció en el capítulo de Alcance, no se contempla dentro del alcance del proyecto un despliegue en un entorno de producción real de terceros como criterio de aceptación del sistema; esta es una restricción reconocida explícitamente, no un límite que el proyecto haya dejado sin identificar.

Finalmente, las cifras del componente Random Forest/Isolation Forest reportadas en el Cuadro \ref{tabla:trazabilidad} (F1 macro, F1 por clase, recall y tasa de falsos positivos de Isolation Forest) provienen de una corrección de trazabilidad de artefactos descrita en la sección de Recolección, procesamiento y validación de datos, integrada a la rama principal del repositorio del proyecto mediante el \textit{pull request} \#79. El análisis SHAP se regeneró posteriormente contra el modelo corregido (\textit{pull request} \#80), de modo que las cifras y los hallazgos de interpretabilidad de este capítulo corresponden al mismo artefacto.
